\documentclass[10pt,letterpaper]{article}
\usepackage{cornell-notes}

\title{Annex A.06: 06-correspondences}
\author{}
\date{}
\setDocTitle{Annex A.06: 06-correspondences}
\setDocSubtitle{ISO/IEC/IEEE 42010:2022}
\setDocOwner{ISO 42010 Cornell Notes}
\setDocDate{\today}
\setCornellCollection{ISO/IEC/IEEE 42010:2022 Cornell Notes}
\setCornellUnitType{Annex}
\setCornellUnitNumber{A.06}
\setCornellUnitTitle{06-correspondences}
\hypersetup{pdftitle={Annex A.06: 06-correspondences},pdfsubject={Cornell notes},pdfauthor={}}

\begin{document}
\maketitle
\begin{CornellOverviewBox}{Overview}
{\Large\bfseries\textcolor{CNavy}{Annex A.6: Correspondences}}\\[3pt]
{\small\textbf{Classification:} Informative annex}\\
{\small\textbf{Learning objective:} Use the informative guidance on correspondences to reinforce the normative and conceptual material without treating the guidance itself as mandatory.}
\end{CornellOverviewBox}
\vspace{3pt}

\begin{CornellNotesTable}
\CornellNoteRow{Purpose / key concept}{Correspondences are used to identify or express named relations within and between AD elements.}
\CornellNoteRow{Core details}{\begin{itemize}[leftmargin=*,nosep,topsep=1pt]
\item In such cases, correspondences can be used to identify and express these named relations.
\item The 2011 edition of This section introduced correspondences and correspondence rules to express and enforce relations between AD elements.
\item Here, AD element is an identified or named part of an architecture description.
\item AD elements include stakeholders, concerns, stakeholder perspectives, and aspects identified in an AD, and views, view components, viewpoints, and model kinds included in an AD.
\end{itemize}}
\CornellNoteRow{Distinction}{A correspondence is the relation itself; a correspondence method specifies rules or practices governing such relations.}
\CornellNoteRow{Study relationship / visual insight}{The correspondence examples show that relations may connect elements inside or across views and ADs, may be governed by methods, and can be n-ary rather than merely binary.}
\CornellNoteRow{Example / application}{Consider two view components in an automotive system’s AD: a software application view component and an electronic control unit (ECU) view component.}
\CornellNoteRow{Related sections}{6.9.1, 6.9.2, 5.2.11, 1.5, 6.8}
\CornellNoteRow{Active recall}{\begin{itemize}[leftmargin=*,nosep,topsep=1pt]
\item What is the central role of Correspondences?
\item How does Correspondences connect to adjacent architecture-description concepts?
\item What closely related concepts must be kept distinct?
\end{itemize}}
\end{CornellNotesTable}


\begin{CornellSummaryBox}{Summary}
Correspondences are used to identify or express named relations within and between AD elements. A correspondence is the relation itself; a correspondence method specifies rules or practices governing such relations. This material is informative guidance and does not itself create a conformance requirement.
\end{CornellSummaryBox}

\end{document}
